% ============================================================
% Chapter 2: State of the Art
%   Section: Automated Fact-Checking
%     Subsection: Closed- and Open-Retrieval Verification
% Included from chapters/02-state-of-the-art/01-automated-fact-checking/automated-fact-checking.tex
% ============================================================

\subsection{Closed- and Open-Retrieval Verification}

Verification systems differ first in where their evidence comes from.

In the \textbf{closed} setting, the evidence is a fixed corpus known in
advance. FEVER, the most influential dataset of the field, built its
claims by altering sentences of Wikipedia and defines the correct
evidence as sentences of the same Wikipedia dump
\cite{thorne2018fever}. Systems index that corpus and are evaluated on
whether they find the annotated sentences and reach the annotated
label. The setting makes evaluation exact and reproducible, and it is
unlike a real fact-check, where nobody knows in advance which document
holds the answer.

In the \textbf{open} setting, evidence is sought on the web. MultiFC
and X-Fact attach to each claim the snippets a search engine returned
for it when the dataset was built
\cite{augenstein2019multifc,gupta2021xfact}; AVeriTeC goes further and
has annotators answer questions about each real-world claim from web
documents, restricted to those published before the claim, so that the
evidence is what a fact-checker could have found at the time
\cite{schlichtkrull2023averitec}.

Open retrieval brings a problem the closed setting does not have:
\emph{leaked evidence}. Glockner, Hou and Gurevych showed that much of
the evidence in such datasets was published after the claim, often by
the fact-checkers themselves, so a system can ``verify'' a claim by
reading its verdict; with only the evidence available before the claim,
performance drops sharply \cite{glockner2022missing}. This matters for
any system that searches today's web for yesterday's claims, as this
one does, and the evaluation flags such cases rather than counting them
(Sections~\ref{sec:rq-metrics} and~\ref{subsec:pilot-time}).

The system built here works in the open setting by necessity: the
claims of constructive news are recent and specific, and no corpus
holds their evidence. Why that rules out FEVER as its benchmark is
argued in Section~\ref{subsec:discard-justification}.
